\chapter{고양이를 알아보는 법}
% full version: chapters/12_recognition.tex

\pencilfig{12}{\pencilart{12}}{큰 고양이와 작은 고양이는 눈의 감지기에 맺히는 그림이 전혀 다르지만, 답은 하나다.}

그림 속 고양이는 왼쪽에 있을 수도 오른쪽에 있을 수도 있고, 가까이 있을 수도 멀리 있을 수도 있으며, 햇빛 아래 있을 수도 그늘에 있을 수도 있다. 그때마다 눈의 감지기에는 전혀 다른 무늬가 떨어지는데도 여러분은 순식간에 고양이를 알아본다. 인식의 어려움이 바로 여기에 있다. 같은 것을 수천 가지 다른 모습에서 알아봐야 한다.

\section*{컨베이어}

뇌는 이 문제를 열 단계 남짓의 컨베이어로 푼다. 신호가 컨베이어를 다 지나는 데 0.15초가 걸린다. 단계마다 두 가지 일이 일어난다. 먼저 “그리고”다. 필요한 조각이 필요한 배치로 놓여 있으면 뉴런이 반응한다. 예를 들어 두 선분이 만나 모서리를 이룰 때다. 다음은 “모든 자리에 걸친 또는”이다. 다음 뉴런은 근처 어디든 모서리가 있으면 반응하고, 정확한 자리는 따지지 않는다. 첫 번째 연산은 인식을 까다롭게 만들고, 두 번째 연산은 위치 이동에 너그럽게 만든다. 단계를 거듭하며 조각은 부분으로, 부분은 물체로 모이고, 위치와 크기의 차이는 용서된다.

\pencilfig{112}{%
\foreach \i/\y/\cx/\cy in {1/1.2/0.15/0.3,2/0.3/0.3/0.1,3/-0.6/0.05/0.05}{
  \draw[penlite] (0,\y-0.3) rectangle (0.6,\y+0.3);
  \draw[pcGray, line width=0.8pt] (\cx,\y-0.3+\cy+0.35) -- (\cx,\y-0.3+\cy) -- (\cx+0.3,\y-0.3+\cy);
  \draw[pen=pcBlue, wash=pcBlue] (1.9,\y) circle (0.2);
  \draw[arr] (0.65,\y) -- (1.65,\y);
  \draw[arr=pcBlue] (2.12,\y) -- (3.62,{0.3+0.12*(2-\i)});}
\draw[pen=pcBlue, wash=pcBlue] (3.9,0.3) circle (0.25);
\draw[arr] (4.2,0.3) -- (5.75,0.3); \node[lbl, anchor=south] at (4.97,0.35) {다음 단계들};
% a small cat
\draw[penlite] (6.3,0.3) circle (0.32);
\draw[penlite] (6.03,0.5) -- (6.07,0.82) -- (6.23,0.6); \draw[penlite] (6.57,0.5) -- (6.53,0.82) -- (6.37,0.6);
\node[lbl, anchor=north] at (0.3,-1.0) {그림}; \node[lbl, anchor=north, align=center] at (1.9,-1.0) {그리고: 이 자리의\\모서리};
\node[lbl, anchor=north, align=center] at (3.9,-1.0) {또는: 어디든\\모서리}; \node[lbl, anchor=north] at (6.3,-1.0) {고양이};
}{컨베이어의 한 단계. “그리고”는 한 자리에서 모서리를 조립하고, “또는”은 그것이 정확히 어디 있는지를 용서한다.}

휴대폰에서 그림을 인식하는 합성곱 신경망이 베낀 것이 바로 이 구조다. 그리고 신경망도 보답을 한다. 물체를 알아보도록 훈련된 신경망은 시각의 윗단계 뉴런이 어떻게 반응할지를 꽤 잘 예측한다.

\section*{동영상에서 배우기}

그런데 뇌는 누가 가르쳤을까? “고양이”라는 이름표가 붙은 그림 백만 장을 여러분에게 보여 준 사람은 없다. 그럴듯한 답은 동영상이다. 세상은 끊김 없이 흘러가고, 1초도 안 되는 사이에 물체 자체는 거의 변하지 않는다. 변하는 것은 위치와 조명이다. 지금 보이는 것과 한순간 전에 보였던 것 사이의 연결이 강해진다면, 네트워크는 한 물체의 여러 모습을 스스로 묶어 낸다. 느리게 변하는 것은 물체이고, 빠르게 변하는 것은 주변 사정이다. 눈이 움직이는 동안 물체를 몰래 바꿔치기한 실험이 있는데, 그러자 뉴런이 서로 다른 물체를 하나로 여기기 시작했다. 일반 규칙으로서는 아직 가설이다.

\section*{착시는 메커니즘의 지문이다}

착시는 고장이 아니라 기계가 어떻게 만들어졌는지를 보여 주는 흔적이다. 폭포를 오래 보다가 가만히 있는 바위를 보면 바위가 위로 “기어오른다”. “아래”와 “위” 채널 한 쌍 가운데 한쪽이 지쳐서 균형이 깨진 것이다. 어떤 사람에게는 흰색과 금색으로, 다른 사람에게는 파란색과 검은색으로 보이는 유명한 드레스는, 각자의 뇌가 말없이 어떤 조명을 가정하느냐에 따라 사람들을 가른다. 그려 놓은 “회전하는 뱀”이 움직이는 것처럼 보이는 까닭은 밝은 부분이 흐린 부분보다 조금 빨리 처리되고, 그 시간차가 움직임으로 읽히기 때문이다. 튀어나왔다가 들어갔다가 하는 정육면체는 서로 경쟁하는 두 뉴런 무리다. 한 모델에 따르면 어느 쪽이 이길지는 주로 잡음이 정한다.

흥미롭게도 동영상의 다음 프레임을 예측하도록 훈련된 신경망도 뱀이 도는 것을 “본다”. 그렇다면 어떤 착시는 예측하는 능력의 피할 수 없는 대가다.

\section*{뇌 대 신경망}

둘이 완전히 닮은 것은 아니다. 뇌는 예시 몇 개로 충분하고, 이름표가 거의 없이 배우며, 모든 일에 20와트를 쓴다. 신경망은 사람 눈에 띄지 않게 조작한 그림에 속을 수 있다. 다만 그런 조작 그림을 사람에게 아주 잠깐 보여 주면 사람도 조금 헷갈린다. 두 기계 사이의 진짜 경계가 어디인지는 아직 밝혀지는 중이다.

\fullversion{12}
